% CASCADE A0 landscape poster (Beamer) with the red side panel
% Compile with:  latexmk -pdf cascade-poster-landscape.tex
%   (or -xelatex / -lualatex to use the Rubik heading font, if installed)
\documentclass[final]{beamer}
\usepackage[size=a0,orientation=landscape,scale=1.4]{beamerposter}
\usetheme[panel]{Cascade}
\usepackage{pgfplots}
\pgfplotsset{compat=1.17}

\title{Your poster title, big and bold}
\author{Author One, Author Two}
\institute{Department of Computer Science and Technology, University of Cambridge}
\event{Conference name \textperiodcentered{} City \textperiodcentered{} Date}
\contact{name@cl.cam.ac.uk}
\qrurl{https://www.cl.cam.ac.uk/research/cascade/}   % the QR code links here
\takeaway{The single result or idea you want people to leave with --- written for a non-specialist.}

\begin{document}
\begin{frame}[t]
\begin{columns}[t,totalwidth=\linewidth]

% ---------------------------------------------------------------- Column 1
\begin{column}{0.315\linewidth}
  \begin{block}{The problem}
    What limits today's processors, accelerators, interconnects or design
    tools --- and why it matters now.
    \begin{itemize}
      \item Constraint or trend one
      \item Constraint or trend two
    \end{itemize}
  \end{block}

  \begin{block}{Our approach}
    Describe the key idea and how it differs from prior work.
    \begin{figure}
      \cascadeplaceholder{46cm}
      \caption{Method or tool-flow diagram.}
    \end{figure}
  \end{block}
\end{column}

% ---------------------------------------------------------------- Column 2
\begin{column}{0.315\linewidth}
  \begin{block}{Design}
    \begin{figure}
      \cascadeplaceholder{46cm}
      \caption{Architecture / layout / die photo.}
    \end{figure}
    Two or three lines walking the reader through the figure.
  \end{block}

  % Headline results: red card with peach numbers
  \setbeamercolor{cascade stat}{fg=cascadepeach}
  \setbeamercolor{cascade stat label}{fg=white}
  \begin{tcolorbox}[enhanced,colback=cascadered,colframe=cascadered,arc=0.6em,boxrule=0pt,
      left=1.4em,right=1.4em,top=1.6em,bottom=3em]
    {\usebeamerfont{block title}\color{white}Headline results}\par\vskip1.5em
    \begin{columns}[t,totalwidth=\linewidth]
      \column{0.32\linewidth}\cascadestat{2.3$\times$}{throughput}
      \column{0.32\linewidth}\cascadestat{$-$41\%}{energy / op}
      \column{0.32\linewidth}\cascadestat{$<$1\%}{area overhead}
    \end{columns}
  \end{tcolorbox}
\end{column}

% ---------------------------------------------------------------- Column 3
\begin{column}{0.315\linewidth}
  \begin{block}{Evaluation}
    \centering
    \begin{tikzpicture}
      \begin{axis}[
          width=\linewidth, height=45cm,
          cascade bars, bar width=1.3cm, ymax=3.6, ytick={0,1,2,3}, yticklabel style={/pgf/number format/precision=0},
          symbolic x coords={Bench A,Bench B,Bench C,Bench D}, xtick=data,
          ylabel={Speed-up over baseline},
          ]
        \addplot[fill=cascadelightgrey] coordinates {(Bench A,1.0) (Bench B,1.0) (Bench C,1.0) (Bench D,1.0)};
        \addplot[fill=cascadered] coordinates {(Bench A,1.6) (Bench B,2.3) (Bench C,1.4) (Bench D,3.1)};
        \legend{Baseline, Our design}
      \end{axis}
    \end{tikzpicture}
  \end{block}

  \begin{block}{Conclusions}
    \begin{itemize}
      \item Main finding, stated plainly
      \item What it enables next
      \item Open questions and future work
    \end{itemize}
    \vskip0.8em
    Where the work goes next, and how industry partners or students can get involved.
    \vskip1.5em
    {\footnotesize\color{cascadegrey}
     \textbf{\color{cascadetext}References}\par
     [1] A.~Author et al., Paper title, \emph{Venue}, Year.\par}
  \end{block}
\end{column}

\end{columns}
\end{frame}
\end{document}
